\SourceLeaf{406}%
\par\medskip\Needspace{6\baselineskip}
\pdfbookmark[2]{VII.}{carnot-zusatz-vii}
\centerline{\textbf{VII.}}\nobreak\smallskip
\begin{center}\textit{Entwickelung der Grundformeln der sphärischen Trigonometrie.}\end{center}
\hspace*{\parindent}Sind \(a,b,c\) die Seiten: \(A,B,C\) die gegenüberstehenden Winkel eines sphärischen Dreiecks, so ist
\[
\cos a=\cos b.\cos c+\sin b.\sin c.\cos A
\]
wie \textsc{Lagrange} für den Fall, dass sowohl \(b\) als \(c\) kleiner wie \(90^0\), elegant bewiesen hat. Indessen lässt sich der Beweis leicht auf alle andere Fälle ausdehnen.

Es können folgende Fälle eintreten:
\[
\text{I.}\quad .\ .\ .\ .\ .\ .\qquad b<90^0,\quad c<90^0
\]
\par
Hier gilt der Beweis unmittelbar.
\[
\text{II.}\quad .\ .\ .\ .\ .\ .\qquad b>90^0,\quad c>90^0
\]
\par
Man verlängere die Seiten \(b\) und \(c\) über die Punkte \(B\) und \(C\) hinaus bis zum Durchschnitte \(A'\), und bestimme den Werth von \(a\) aus der Betrachtung des Dreiecks \(A'BC\).
\[
\text{III.}\quad .\ .\ .\ .\ .\ .\qquad b>90^0,\quad c<90^0
\]
\par
Man verlängere die Seiten \(b\) und \(a\) über die Punkte \(A\) und \(B\) hinaus bis zum Durchschnitte \(C'\), in dem Dreieck \(C'AB\) ist sodann aus Fall I
\[
\cos(180^0-a)=\cos c.\cos(180^0-b)+\sin c.\sin(180^0-b).\cos(180^0-A)
\]
welches mit der Grundformel identisch ist.

Mit diesem Falle ist \(b<90^0\) und \(c>90^0\) wesentlich einerlei.
\[
\text{IV.}\quad .\ .\ .\ .\ .\ .\qquad b=90^0,\quad A=90^0
\]
\par
Hier ist \(C\) der Pol von \(AB\), also nothwendig auch \(a=90^0\). Folglich ist die Formel
\[
\cos a=\cos b.\cos c+\sin b.\sin c.\cos A
\]
von selbst evident. Mit diesem Falle ist \(c=90^0,\ A=90^0\) wesentlich einerlei.
